\clearpage
\hypertarget{native-484-psyche}{}
\begin{center}
\rule{0.55\linewidth}{0.4pt}\\[2pt]
Astronomische Nachrichten. Vol. XXXV. No. 818. p. 17. 27 July 1852.\\[2pt]
\rule{0.55\linewidth}{0.4pt}
\end{center}
\begin{center}Elements of \emph{Psyche}, with remarks on determining the orbit of this planet.\end{center}

Because the inclination of Psyche’s orbit is so very small, three observations are quite unsuitable for determining its orbit precisely. The elements that several people obtained this way confirm this as well. I therefore arranged for Mr \textsc{Klinkerfues} to base a new determination on four positions, using the method developed in the Th. M. C. C., which seems to have been used little, if at all, so far. I enclose his elements here, together with a comparison with all the observations that have come to our attention.

The close agreement with \textsc{Hind}’s estimate of 29 January is interesting. It leaves no room for further doubt that the star \textsc{Hind} saw then is the same as the planet discovered by \textsc{Gasparis} 7 weeks later. If Mr \textsc{Hind} were willing to give more details about his record at the time, especially about the hour and how accurate he considers his estimate, including this observation in a new calculation could probably make the result appreciably more accurate still.

\Needspace{17\baselineskip}
\begin{center}
Elements of Psyche.\\
Epoch: March 31.0, Berlin mean time.
\end{center}
\[
\text{Mean longitude}=151^\circ\ 32'\ 5''56
\]
\[
\left.\begin{array}{rcl}
\pi &=& 6\ 20\ 2.86 \\
\AscendingNode &=& 150\ 33\ 5.86
\end{array}\right\}\quad\text{Mean equinox 1852.0}
\]
\[
\begin{array}{rcl}
i &=& 3\ 2\ 16.06 \\
\varphi &=& 6\ 38\ 32.80 \\
\log e &=& 9.0632320 \\
\log a &=& 0.4693195 \\
\mu &=& 701''4993
\end{array}
\]

The elements above were calculated from 4 normal positions for 19 March, 14 April, 30 April and 14 May.

Among the observations compared below [not reproduced here] with this system is Mr \textsc{Hind}’s estimate of 29 January. Since the time of that estimate is unknown, it was assumed to correspond to about midnight in Berlin. At that time the planet’s hourly motion is $22''75$ westward and $10''25$ northward.

